\documentclass[11pt,a4paper]{article}

\usepackage{fontspec}
\setmainfont{DejaVu Sans}[Scale=0.90]
\usepackage{amsmath,amssymb}
\usepackage[margin=2.4cm]{geometry}
\usepackage{booktabs,array,longtable,tabularx}
\usepackage{enumitem}
\usepackage{titlesec}
\usepackage{parskip}

\titlelabel{\thetitle.\quad}
\titleformat{\section}{\normalfont\large\bfseries}{\thesection.}{0.6em}{}
\titleformat{\subsection}{\normalfont\normalsize\bfseries}{\thesubsection}{0.6em}{}
\titlespacing*{\section}{0pt}{16pt}{6pt}
\titlespacing*{\subsection}{0pt}{12pt}{4pt}

\setlist[itemize]{leftmargin=1.4em,itemsep=1pt,topsep=3pt}
\newcolumntype{L}[1]{>{\raggedright\arraybackslash}p{#1}}

\begin{document}

\begin{center}
{\Large\bfseries Corrections, démonstrations et rappels de mécanique quantique}\\[4pt]
{\normalsize Document d'accompagnement des notes 01 à 04 sur Krapez (2018)}\\[2pt]
{\small 8 septembre 2026}
\end{center}

\vspace{6pt}

Ce document comporte quatre parties. La première recense les erreurs relevées dans les fichiers 01 à 04
et donne pour chacune la vérification qui tranche. La deuxième contient les onze démonstrations écrites
en entier, sans étape sautée. La troisième expose ce qu'est réellement l'équation de Schrödinger, afin
que l'analogie puisse être décrite sans erreur. La quatrième liste les formulations à ne pas employer.

\section{Errata}

\subsection{Signe du potentiel : inversé dans deux fichiers sur trois}

Le fichier 01 (§1.5) et le fichier 02 (C25) affirment que le potentiel négatif donne les solutions
hyperboliques et le potentiel positif les solutions trigonométriques. Le fichier 03 affirme le
contraire. Les trois documents ne peuvent pas avoir raison ensemble.

\textbf{Le fichier 03 a raison.} Krapez le fixe sans ambiguïté à son équation (32) : cosinus et sinus
hyperboliques si $\beta > 0$, cosinus et sinus circulaires si $\beta < 0$.

\textbf{La vérification.} Le potentiel est défini par $V = s''/s$. Imposer $V = \beta$ revient donc à
résoudre
\[ s'' = \beta\, s . \]

Si $\beta > 0$, la dérivée seconde a le \emph{même} signe que la fonction : la courbure éloigne la
courbe de l'axe et la solution s'emballe. Ce sont les exponentielles, donc $\cosh$ et $\sinh$.

Si $\beta < 0$, la dérivée seconde a le signe \emph{opposé} : la courbure ramène toujours la courbe
vers l'axe. C'est un ressort, donc une oscillation, donc $\cos$ et $\sin$.

Moyen mnémotechnique : $\beta$ négatif $=$ rappel $=$ ressort $=$ oscillation ; $\beta$ positif $=$
répulsion $=$ divergence exponentielle.

\subsection{Relations inverses entre propriétés : racines manquantes}

Le fichier 03 (§2.1 et §3.2) écrit $\lambda = a\,b$ et $\rho c = b/a$. Les deux sont faux : il manque
les racines carrées. Les relations correctes sont
\[ \lambda = b\sqrt{a}
\qquad\text{et}\qquad
\rho c = \frac{b}{\sqrt{a}} . \]

La démonstration est en §2.1 ci-dessous. C'est l'erreur la plus grave du lot : ces deux identités sont
exactement ce qui fait se simplifier les racines à la première étape de Liouville. Avec la version
fausse, toute la transformation s'écroule.

\subsection{Référence bibliographique non remplie}

Le fichier 01 contient encore un espace réservé : « [Référence complète conservée selon le document
source] ». Sur une fiche dont la raison d'être est la vérification des références, c'est le champ à ne
pas laisser vide. Le titre exact est : \emph{Linear, trigonometric and hyperbolic profiles of thermal
effusivity in the Liouville space and related quadrupoles: Simple analytical tools for modeling graded
layers and multilayers.}

\subsection{Famille à exposant libre : formule inexacte}

Le fichier 02 (C22) écrit une relation de la forme $z_q = \int (A\,a^q)\,d\xi$. Ce n'est pas la
relation de l'article. La bonne est
\[ z(\xi) \;=\; \int_0^{\xi} \lambda^{q}\, c^{\,q-1}\, b^{\,1-2q} \, du , \]
valable pour \emph{toute} valeur de $q$. La démonstration §2.8 le vérifie. Le tableau des trois valeurs
particulières, lui, est correct dans le fichier.

\subsection{Signe du flux dans le vecteur d'état}

Le fichier 02 (C17) écrit le couple transformé $(T,\; b\,T_\xi)$. Le flux vaut $-\,b\,\theta'$, avec un
signe moins. Le couple correct est $(\theta,\; -b\,\theta')$. Sans ce signe, les entrées du quadripôle
changent de signe et le test de la couche homogène échoue.

\subsection{Le potentiel n'a cette forme que dans un des deux cas}

Le fichier 02 (C12) et le fichier 03 (§3.5) donnent
\[ V(\xi) = \frac{1}{\sqrt{b}}\,\frac{d^{2}\sqrt{b}}{d\xi^{2}} . \]
C'est juste, mais seulement pour la formulation en température. La forme générale est $V = s''/s$ avec
$s = b^{\pm 1/2}$, l'exposant positif pour la température et le négatif pour le flux. Comme les
fichiers insistent ailleurs sur la dualité des deux formulations, il faut être cohérent ici aussi.

\newpage
\section{Les onze démonstrations}

Chacune est courte et aucune étape n'est sautée. Il faut les refaire au stylo : c'est la seule façon
de pouvoir les défendre au tableau.

\subsection{Les deux identités du changement de coordonnée}

\textbf{À démontrer :} $\lambda = b\sqrt{a}$ et $\rho c = b/\sqrt{a}$, à partir des seules définitions
$a = \lambda/\rho c$ et $b = \sqrt{\lambda \rho c}$.

\textbf{Démonstration.}
\[ b\cdot\sqrt{a} \;=\; \sqrt{\lambda\rho c}\cdot\sqrt{\frac{\lambda}{\rho c}}
\;=\; \sqrt{\lambda\rho c \cdot \frac{\lambda}{\rho c}} \;=\; \sqrt{\lambda^{2}} \;=\; \lambda \]
\[ \frac{b}{\sqrt{a}} \;=\; \frac{\sqrt{\lambda\rho c}}{\sqrt{\lambda/\rho c}}
\;=\; \sqrt{\lambda\rho c \cdot \frac{\rho c}{\lambda}} \;=\; \sqrt{(\rho c)^{2}} \;=\; \rho c \]

\textbf{Corollaires utilisés partout dans l'article :}
\[ \frac{\lambda}{\sqrt{a}} = b
\qquad\text{et}\qquad
\frac{1}{\rho c\,\sqrt{a}} = \frac{1}{b} . \]
Ce sont ces deux simplifications qui font disparaître la diffusivité à l'étape 1.

\subsection{Analyse dimensionnelle complète}

\textbf{Diffusivité.}
$[a] = \dfrac{\mathrm{W\,m^{-1}K^{-1}}}{\mathrm{J\,m^{-3}K^{-1}}} = \mathrm{m^{2}\,s^{-1}}$,
puisqu'un watt est un joule par seconde.

\textbf{Effusivité.}
$[b] = \sqrt{\mathrm{W\,m^{-1}K^{-1}} \cdot \mathrm{J\,m^{-3}K^{-1}}}
= \mathrm{W\,s^{1/2}\,m^{-2}\,K^{-1}}$.

\textbf{Coordonnée transformée.}
$[\xi] = \dfrac{\mathrm{m}}{\mathrm{m\,s^{-1/2}}} = \mathrm{s^{1/2}}$.
Son carré est donc un temps : le temps de diffusion.

\textbf{Potentiel.} $V = s''/s$. Deux dérivations par rapport à $\xi$ divisent par
$(\mathrm{s^{1/2}})^{2} = \mathrm{s}$, et le rapport à $s$ élimine les unités de $s$. Donc
$[V] = \mathrm{s^{-1}}$.

\textbf{Paramètre spectral.} $p = i\omega$ est en $\mathrm{rad\,s^{-1}}$, donc en $\mathrm{s^{-1}}$.

\textbf{Vérification de cohérence.} $V$ et $p$ ont la même dimension, ce qui est indispensable
puisqu'ils s'additionnent dans l'équation. Si l'on trouve autre chose, une étape est fausse.

\subsection{Le changement de coordonnée, en entier}

\textbf{Point de départ} --- l'équation transformée en temps, à propriétés variables :
\[ p\,\rho c(z)\,\theta \;=\; \frac{d}{dz}\!\left[\lambda(z)\,\frac{d\theta}{dz}\right] . \]

\textbf{Définition de la nouvelle coordonnée :}
\[ \xi(z) \;=\; \int_{0}^{z} \frac{du}{\sqrt{a(u)}} . \]

\textbf{Règle de la chaîne.} Dériver la définition donne $d\xi/dz = 1/\sqrt{a}$, donc
\[ \frac{d}{dz} \;=\; \frac{d\xi}{dz}\,\frac{d}{d\xi} \;=\; \frac{1}{\sqrt{a}}\,\frac{d}{d\xi} . \]

\textbf{Membre de gauche.} On applique la règle deux fois. À l'intérieur d'abord :
\[ \lambda\,\frac{d\theta}{dz}
\;=\; \lambda\cdot\frac{1}{\sqrt{a}}\cdot\frac{d\theta}{d\xi}
\;=\; \frac{\lambda}{\sqrt{a}}\,\theta' \;=\; b\,\theta' , \]
en utilisant le corollaire de §2.1. Puis à l'extérieur :
\[ \frac{d}{dz}\bigl[\,b\,\theta'\,\bigr]
\;=\; \frac{1}{\sqrt{a}}\,\frac{d}{d\xi}\bigl[\,b\,\theta'\,\bigr]
\;=\; \frac{1}{\sqrt{a}}\,\bigl(b\,\theta'\bigr)' . \]

\textbf{Membre de droite.} Avec $\rho c = b/\sqrt{a}$, également de §2.1 :
\[ p\,\rho c\,\theta \;=\; p\,\frac{b}{\sqrt{a}}\,\theta . \]

\textbf{Assemblage.} Les deux membres portent le même facteur $1/\sqrt{a}$ : on multiplie tout par
$\sqrt{a}$ et il reste
\[ \boxed{\; p\,b(\xi)\,\theta \;=\; \bigl(b(\xi)\,\theta'\bigr)' \;} \]

\textbf{Ce qu'il faut voir.} La diffusivité n'apparaît plus. Deux fonctions matériau inconnues sont
devenues une seule, l'effusivité. Le même calcul sur la formulation en flux donne l'équation duale,
avec $b$ remplacée par $1/b$.

\newpage
\subsection{Pourquoi la racine de l'effusivité est le seul facteur possible}

\textbf{Point de départ.} On développe le résultat précédent, $b\,\theta'' + b'\,\theta' = p\,b\,\theta$,
puis on divise par $b$ :
\[ \theta'' + \frac{b'}{b}\,\theta' \;=\; p\,\theta . \]
Le terme du milieu, en dérivée première, est celui qu'on veut supprimer.

\textbf{La méthode.} On ne suppose pas la réponse. On pose $\theta = \psi\,u$ avec $u$ \emph{inconnue},
on développe, et on impose que le coefficient de $\psi'$ s'annule. Cela déterminera $u$.
\[ \theta' = \psi'u + \psi u'
\qquad
\theta'' = \psi''u + 2\psi'u' + \psi u'' \]

En reportant :
\[ \psi''u + 2\psi'u' + \psi u'' + \frac{b'}{b}\bigl(\psi'u + \psi u'\bigr) \;=\; p\,\psi\,u . \]

\textbf{On isole le coefficient de $\psi'$ et on l'annule :}
\[ 2u' + \frac{b'}{b}\,u \;=\; 0 . \]

\textbf{On résout cette équation en $u$}, qui est à variables séparables :
\[ 2\,\frac{u'}{u} = -\frac{b'}{b}
\qquad\Longrightarrow\qquad
2\ln|u| = -\ln|b| + \text{const}
\qquad\Longrightarrow\qquad
u = C\,b^{-1/2} . \]

La constante $C$ ne change rien à l'équation finale ; on la prend égale à 1. Donc $u = 1/\sqrt{b}$,
c'est-à-dire
\[ \boxed{\; \psi = \sqrt{b}\;\theta \;} \]

\textbf{Conclusion.} Le facteur n'a pas été choisi, il a été \emph{calculé}. Aucune autre fonction
n'annule ce terme, puisque l'équation en $u$ est du premier ordre et n'a qu'une famille de solutions.

\subsection{La forme du potentiel}

On reprend le calcul précédent avec $u = b^{-1/2}$. Le terme en $\psi'$ ayant disparu, il reste
\[ \psi''u + \psi\left[\,u'' + \frac{b'}{b}\,u'\,\right] \;=\; p\,\psi\,u , \]
puis, en divisant par $u$,
\[ \psi'' + \psi\left[\,\frac{u''}{u} + \frac{b'}{b}\cdot\frac{u'}{u}\,\right] \;=\; p\,\psi . \]

\textbf{Calcul des deux rapports.} Avec $u = b^{-1/2}$ :
\[ u' = -\tfrac{1}{2}\,b^{-3/2}b'
\qquad\Longrightarrow\qquad
\frac{u'}{u} = -\frac{1}{2}\,\frac{b'}{b} \]
\[ u'' = -\tfrac{1}{2}\,b^{-3/2}b'' + \tfrac{3}{4}\,b^{-5/2}(b')^{2}
\qquad\Longrightarrow\qquad
\frac{u''}{u} = -\frac{1}{2}\,\frac{b''}{b} + \frac{3}{4}\left(\frac{b'}{b}\right)^{2} \]

\textbf{Somme des deux termes entre crochets :}
\[ -\frac{1}{2}\frac{b''}{b} + \frac{3}{4}\left(\frac{b'}{b}\right)^{2}
- \frac{1}{2}\left(\frac{b'}{b}\right)^{2}
= -\frac{1}{2}\frac{b''}{b} + \frac{1}{4}\left(\frac{b'}{b}\right)^{2} . \]

\textbf{Reconnaissance.} Posons $s = \sqrt{b}$ et calculons $s''/s$ :
\[ s' = \tfrac{1}{2}\,b^{-1/2}b'
\qquad
s'' = \tfrac{1}{2}\,b^{-1/2}b'' - \tfrac{1}{4}\,b^{-3/2}(b')^{2} \]
\[ \frac{s''}{s} = \frac{1}{2}\frac{b''}{b} - \frac{1}{4}\left(\frac{b'}{b}\right)^{2} \]

C'est l'opposé exact du crochet. L'équation devient donc
\[ \boxed{\;\frac{d^{2}\psi}{d\xi^{2}} - \bigl(V(\xi) + p\bigr)\psi = 0
\qquad\text{avec}\qquad
V(\xi) = \frac{s''(\xi)}{s(\xi)},\quad s = b^{\pm 1/2} \;} \]

\textbf{Lecture physique du potentiel.} C'est la courbure relative de la métapropriété. Là où
l'effusivité varie de sorte que sa racine est linéaire en $\xi$, la courbure est nulle et le potentiel
s'annule.

\newpage
\subsection{La loi de Fourier dans l'espace transformé}

\textbf{Point de départ :} la loi de Fourier ordinaire, $\phi = -\lambda\, d\theta/dz$.

\textbf{On applique la règle de la chaîne de §2.3 :}
\[ \phi \;=\; -\lambda\cdot\frac{1}{\sqrt{a}}\cdot\frac{d\theta}{d\xi}
\;=\; -\frac{\lambda}{\sqrt{a}}\,\theta'
\;=\; \boxed{\;-\,b(\xi)\,\theta'\;} \]

\textbf{Conséquence.} La conductivité a été remplacée par l'effusivité. Toute condition aux limites
portant sur un flux --- un laser qui chauffe, une face adiabatique --- s'écrira donc dans l'espace
$\xi$ avec $b$ seule.

C'est cette propriété qui referme la démonstration de la prééminence : l'équation ne contient que $b$,
les conditions aux limites ne contiennent que $b$, donc la solution ne dépend que de $b$. Et comme la
surface correspond à $\xi = 0$ pour tout matériau, deux milieux de même $b(\xi)$ donnent une réponse de
surface identique à toute fréquence.

\subsection{Le flux en fonction du wronskien}

\textbf{À démontrer :} $\phi = -W(s,\psi)$, avec $W(f,g) = f g' - f' g$, dans la formulation en
température.

\textbf{Démonstration.} On part de $\theta = \psi/s$ et on dérive :
\[ \theta' = \frac{\psi' s - \psi s'}{s^{2}} = \frac{W(s,\psi)}{s^{2}} . \]
On applique §2.6, $\phi = -b\,\theta'$, en se rappelant que dans le cas température $b = s^{2}$ :
\[ \phi = -s^{2}\cdot\frac{W(s,\psi)}{s^{2}} = -W(s,\psi) . \]

\textbf{Rappel sur le wronskien.} Il mesure l'indépendance de deux fonctions : nul, elles sont
proportionnelles. Pour une équation en forme normale --- sans terme de dérivée première, ce qui est
précisément notre cas après §2.4 --- le théorème d'Abel garantit que le wronskien de deux solutions est
\emph{constant}. C'est de là que vient le déterminant unité du quadripôle.

\subsection{L'identité à exposant libre}

\textbf{À démontrer :} $\lambda^{q}\,c^{\,q-1}\,b^{\,1-2q} = \sqrt{a}$, pour toute valeur de $q$.

\textbf{Démonstration.} On remplace $b$ par $(\lambda c)^{1/2}$ et on regroupe les exposants :
\[ \lambda^{q}\,c^{\,q-1}\,(\lambda c)^{(1-2q)/2}
\;=\; \lambda^{\,q + \frac{1-2q}{2}}\; c^{\,q-1+\frac{1-2q}{2}} . \]

L'exposant de $\lambda$ vaut $q + \tfrac{1}{2} - q = \tfrac{1}{2}$ ; celui de $c$ vaut
$q - 1 + \tfrac{1}{2} - q = -\tfrac{1}{2}$. Donc
\[ = \lambda^{1/2}\,c^{-1/2} = \sqrt{\frac{\lambda}{c}} = \sqrt{a} . \]

\textbf{Ce que cela signifie.} Le paramètre $q$ est libre : il ne change rien à la valeur de
l'intégrale, seulement à son écriture. Son intérêt est que si le préfacteur $\lambda^{q}c^{\,q-1}$ se
trouve constant pour une valeur de $q$, il sort de l'intégrale et il ne reste qu'une puissance de $b$
à intégrer.

\medskip
\begin{center}
\begin{tabular}{@{}cl@{}}
\toprule
$q$ & Propriété supposée constante \\
\midrule
$0$ & capacité thermique volumique \\
$1/2$ & diffusivité \\
$1$ & conductivité \\
\bottomrule
\end{tabular}
\end{center}
\medskip

\textbf{La conséquence, et c'est le cœur du sujet.} Chaque valeur de $q$ correspond à une hypothèse
différente sur le matériau, et produit un profil $b(z)$ différent --- à partir du même $b(\xi)$, donc
à partir de la même mesure.

\newpage
\subsection{Le quadripôle homogène, le test unitaire}

\textbf{Situation.} Couche homogène : $b$ constante, donc $s$ constante, donc $s'' = 0$, donc $V = 0$.
L'équation devient
\[ \frac{d^{2}\psi}{d\xi^{2}} = p\,\psi , \]
dont les solutions indépendantes sont $\cosh(\sqrt{p}\,\xi)$ et $\sinh(\sqrt{p}\,\xi)$.

\textbf{Retour aux grandeurs physiques.} $b$ étant constante, $\theta = \psi/\sqrt{b}$ n'est qu'un
changement d'échelle, et $\phi = -b\,\theta'$. En imposant l'état à la face d'entrée et en propageant
jusqu'à $\xi_{1}$, on obtient
\[ A = D = \cosh\bigl(\sqrt{p}\,\xi_{1}\bigr),
\qquad
B = \frac{\sinh\bigl(\sqrt{p}\,\xi_{1}\bigr)}{b\sqrt{p}},
\qquad
C = b\sqrt{p}\,\sinh\bigl(\sqrt{p}\,\xi_{1}\bigr) . \]

\textbf{Vérification obligatoire :}
\[ AD - BC = \cosh^{2} - \sinh^{2} = 1 . \]

\textbf{Lien avec la profondeur physique.} Pour une couche homogène d'épaisseur $D_{\text{ép}}$, on a
$\xi_{1} = D_{\text{ép}}/\sqrt{a}$, donc $\sqrt{p}\,\xi_{1} = D_{\text{ép}}\sqrt{p/a}$. On retrouve
exactement le quadripôle classique du livre de Maillet.

C'est le premier test à écrire dans \texttt{tests/}, avant toute autre ligne de code. Un quadripôle
gradué qui ne redonne pas cette expression quand on égalise les effusivités aux deux faces est faux.

\subsection{Cattaneo : le décalage spectral et la rotation de l'énergie}

\textbf{Loi de Cattaneo.} Le flux ne répond plus instantanément au gradient, il met un temps $\tau$ :
\[ \tau\,\frac{\partial\varphi}{\partial t} + \varphi = -\lambda\,\frac{\partial T}{\partial z} . \]

\textbf{Passage au domaine spectral.} La dérivée temporelle devient une multiplication par $p$ :
\[ \phi\,(1 + \tau p) = -\lambda\,\frac{d\theta}{dz}
\qquad\Longrightarrow\qquad
\phi = -\frac{\lambda}{1+\tau p}\,\frac{d\theta}{dz} . \]

\textbf{Combinaison avec le bilan d'énergie} $p\,\rho c\,\theta = -\,d\phi/dz$ :
\[ \frac{d}{dz}\!\left[\lambda\,\frac{d\theta}{dz}\right] = p\,(1+\tau p)\,\rho c\,\theta . \]

\textbf{Résultat.} Toute la structure de l'équation est inchangée ; seul le paramètre spectral a été
remplacé :
\[ p \;\longrightarrow\; p_{\text{eff}} = p + \tau p^{2} . \]

\textbf{En régime modulé,} $p = i\omega$, donc $p^{2} = -\omega^{2}$ et
\[ p_{\text{eff}} = i\omega - \tau\omega^{2} . \]

\textbf{L'énergie de l'analogie} vaut $E = -p_{\text{eff}}$ :
\[ \boxed{\; E = \tau\omega^{2} - i\,\omega = \omega\,(x - i)
\qquad\text{avec}\qquad x = \omega\tau \;} \]

\textbf{Son argument.} La partie réelle vaut $\tau\omega^{2}$, la partie imaginaire $-\omega$, donc
\[ \arg(E) = \arctan\!\left(\frac{-\omega}{\tau\omega^{2}}\right)
= -\arctan\!\left(\frac{1}{\omega\tau}\right) . \]

\textbf{Les trois limites.} Quand $\omega\tau \to 0$, l'arc tangente de l'infini vaut $90^{\circ}$,
donc $\arg(E) \to -90^{\circ}$ : l'énergie est sur l'axe imaginaire, c'est exactement le cas de Krapez.
À $\omega\tau = 1$, $\arctan(1) = 45^{\circ}$, donc $\arg(E) = -45^{\circ}$ exactement. Quand
$\omega\tau \to \infty$, $\arg(E) \to 0^{\circ}$ : l'énergie approche l'axe réel positif, celui des
états de diffusion.

\textbf{Le point qui compte.} L'angle ne dépend que du produit $\omega\tau$. Ni du matériau, ni de
l'épaisseur, ni de la géométrie.

\newpage
\subsection{Le demi-cercle de Cole-Cole}

\textbf{Contexte.} Le module de $E$ croît sans borne, ce qui rend le tracé illisible. On trace donc un
multiple réel positif de $E$, qui a le même argument. Le choix $E^{*} = x/(x+i)$ donne une courbe
remarquable.

\textbf{Première étape --- vérifier que l'argument est le même.} On multiplie par le conjugué :
\[ E^{*} = \frac{x}{x+i} = \frac{x(x-i)}{(x+i)(x-i)} = \frac{x^{2}-ix}{x^{2}+1}
= \frac{x}{x^{2}+1}\,(x-i) . \]
Or $E = \omega\,(x-i)$ d'après §2.10. Le rapport entre les deux vaut donc
\[ \frac{E^{*}}{E} = \frac{x}{\omega\,(x^{2}+1)} = \frac{\tau}{1+\omega^{2}\tau^{2}} \;>\; 0 . \]
Un réel strictement positif ne change pas l'argument. Les deux quantités ont donc rigoureusement le
même argument, ce que le calcul numérique confirme à $10^{-14}$ degré près.

\textbf{Deuxième étape --- montrer que c'est un demi-cercle exact.} On calcule la distance de $E^{*}$
au point $1/2$ :
\[ E^{*} - \frac{1}{2}
= \frac{2x^{2} - 2ix - x^{2} - 1}{2\,(x^{2}+1)}
= \frac{x^{2} - 1 - 2ix}{2\,(x^{2}+1)} . \]

Module du numérateur :
\[ \bigl|\,x^{2}-1-2ix\,\bigr|^{2} = (x^{2}-1)^{2} + 4x^{2}
= x^{4} - 2x^{2} + 1 + 4x^{2} = x^{4} + 2x^{2} + 1 = (x^{2}+1)^{2} . \]

Donc
\[ \boxed{\; \left|\,E^{*} - \tfrac{1}{2}\,\right| = \frac{x^{2}+1}{2\,(x^{2}+1)} = \frac{1}{2} \;} \]

\textbf{Conclusion.} Quelle que soit la valeur de $\omega\tau$, le point est à distance exactement
$1/2$ du point $1/2$. C'est un demi-cercle de rayon $1/2$ centré sur $1/2$.

\textbf{Pourquoi cela compte.} Ce demi-cercle porte un nom : c'est le \emph{diagramme de Cole-Cole},
signature d'un processus à temps de relaxation unique. Il est utilisé depuis les années 1940 en
spectroscopie diélectrique et en spectroscopie d'impédance. La communauté qui l'emploie sait aussi lire
les écarts : un arc aplati, plus bas qu'un demi-cercle, signale une \emph{distribution} de temps de
relaxation plutôt qu'un seul. C'est un diagnostic directement transposable à des données d'AlN, et une
piste bibliographique à vérifier.

\newpage
\section{Ce qu'est réellement l'équation de Schrödinger}

Cette partie existe pour une raison précise : pouvoir parler de l'analogie sans dire de bêtises. Elle
décrit la vraie mécanique quantique, puis établit point par point ce que l'analogie autorise.

\subsection{D'où vient l'équation}

En 1924, Louis de Broglie propose que la matière possède un caractère ondulatoire : à une particule de
quantité de mouvement $p$ est associée une longueur d'onde $\lambda = h/p$, où $h$ est la constante de
Planck. L'idée est vérifiée expérimentalement dès 1927 par la diffraction d'électrons sur un cristal.

En 1926, Erwin Schrödinger écrit l'équation d'évolution de cette onde. Elle joue en mécanique quantique
le rôle que la seconde loi de Newton joue en mécanique classique :
\[ i\hbar\,\frac{\partial\Psi}{\partial t} = \hat{H}\,\Psi , \]
où $\Psi$ est la \emph{fonction d'onde}, $\hbar$ la constante de Planck réduite (environ
$1{,}05\times10^{-34}$ J\,s) et $\hat{H}$ l'\emph{opérateur hamiltonien}, qui représente l'énergie
totale.

\subsection{Ce que représente la fonction d'onde}

C'est le point le plus contre-intuitif, et celui qu'il faut avoir compris avant de parler d'analogie.

$\Psi$ n'est \emph{pas} la position de la particule, ni sa trajectoire, ni une densité de matière.
C'est une \emph{amplitude de probabilité}, à valeurs complexes. Son interprétation, due à Max Born la
même année, est que le module au carré donne la densité de probabilité de présence :
\[ |\Psi(x)|^{2}\,dx = \text{probabilité de trouver la particule entre } x \text{ et } x+dx . \]

Conséquence immédiate : la somme des probabilités sur tout l'espace doit valoir 1,
\[ \int_{-\infty}^{+\infty} |\Psi(x)|^{2}\,dx = 1 . \]

On dit que $\Psi$ doit être \emph{de carré intégrable}. C'est cette contrainte qui produit toute la
quantification, et c'est précisément elle qui \textbf{n'existe pas} chez Krapez.

\subsection{La version stationnaire, et un parallèle familier}

Si le potentiel ne dépend pas du temps, on cherche des solutions de la forme
\[ \Psi(x,t) = \psi(x)\,e^{-iEt/\hbar} . \]

\textbf{C'est exactement le geste employé pour l'équation de la chaleur} : séparer une partie spatiale
et une exponentielle temporelle, pour que la dérivée en temps devienne une multiplication. Là-bas la
dérivée donnait un facteur $p$ ; ici elle donne un facteur $-iE/\hbar$.

En reportant et en simplifiant l'exponentielle, il reste l'équation de Schrödinger stationnaire :
\[ \hat{H}\psi = E\psi
\qquad\text{soit}\qquad
-\frac{\hbar^{2}}{2m}\,\frac{d^{2}\psi}{dx^{2}} + V(x)\,\psi = E\,\psi . \]

$E$ est l'énergie de la particule, un nombre réel. $V(x)$ est l'énergie potentielle : le paysage de
collines et de vallées dans lequel la particule évolue.

\subsection{D'où sort le terme en $\hbar^{2}/2m$}

Ce n'est pas un coefficient décoratif : c'est l'\emph{énergie cinétique}.

En mécanique classique, l'énergie cinétique vaut $p^{2}/2m$. En mécanique quantique, les grandeurs
physiques deviennent des opérateurs, et la quantité de mouvement devient
\[ \hat{p} = -i\hbar\,\frac{d}{dx} . \]
En élevant au carré et en divisant par $2m$ :
\[ \frac{\hat{p}^{2}}{2m} = -\frac{\hbar^{2}}{2m}\,\frac{d^{2}}{dx^{2}} . \]

\textbf{Mise à l'échelle.} En divisant toute l'équation par $\hbar^{2}/2m$ et en renommant le potentiel
et l'énergie, on obtient la forme compacte utilisée pour la comparaison :
\[ \frac{d^{2}\psi}{dx^{2}} = \bigl(\tilde{V}(x) - \tilde{E}\bigr)\,\psi . \]
Rien n'a été négligé : les constantes ont été absorbées dans les définitions. C'est un changement
d'unités, comme travailler en kilomètres plutôt qu'en mètres.

\newpage
\subsection{Valeurs propres : pourquoi l'énergie est quantifiée}

Voici le résultat le plus important de toute la mécanique quantique, et il découle de la contrainte
d'intégrabilité.

Pour une valeur de $E$ choisie au hasard, l'équation stationnaire admet toujours des solutions
mathématiques. Mais presque toutes \emph{divergent} à l'infini : leur intégrale de $|\psi|^{2}$ est
infinie, elles ne représentent donc aucune probabilité et sont physiquement inacceptables.

Seules certaines valeurs particulières de $E$ donnent une solution qui reste finie et intégrable. Ces
valeurs s'appellent les \emph{valeurs propres} de l'opérateur, et les solutions correspondantes ses
\emph{fonctions propres}.

C'est l'origine physique de la quantification. Un atome n'a pas un continuum d'énergies possibles mais
des niveaux discrets, et c'est pour cette seule raison : les autres énergies ne produisent pas de
solution acceptable. Les raies spectrales observées sont les différences entre ces niveaux.

\subsection{États liés et états de diffusion}

La distinction se fait sur la position de $E$ par rapport au potentiel à l'infini. Prenons la
convention usuelle $V \to 0$ loin de la particule.

\textbf{États liés, $E < 0$.} La particule est piégée dans une vallée du potentiel. La fonction d'onde
décroît exponentiellement de part et d'autre : la particule ne peut pas s'échapper. Le spectre est
\emph{discret}. Exemple : l'électron dans un atome.

\textbf{États de diffusion, $E > 0$.} La particule vient de l'infini, rencontre le potentiel, et
repart. Elle n'est pas piégée. La fonction d'onde oscille à l'infini au lieu de décroître, et n'est pas
de carré intégrable au sens strict. Le spectre est \emph{continu} : toutes les énergies positives sont
permises.

Pour ces états on ne parle plus de niveaux mais de coefficients de réflexion et de transmission :
quelle proportion de l'onde rebondit sur l'obstacle, quelle proportion le traverse. C'est la
\emph{théorie de la diffusion}.

\textbf{Résonances.} À certaines énergies, la transmission passe brusquement par un maximum : l'onde
reste temporairement piégée avant de repartir. Mathématiquement, ces résonances sont des pôles situés
dans le plan complexe des énergies, légèrement en dessous de l'axe réel.

\subsection{Le problème inverse de la diffusion}

Voilà le point qui concerne directement le sujet, et le seul emprunt réellement légitime.

La question directe est : connaissant le potentiel $V(x)$, que valent les coefficients de réflexion et
de transmission ? La question inverse est : connaissant ce qui ressort, peut-on reconstruire le
potentiel, de façon unique, et avec quelle stabilité ?

C'est exactement la structure du problème posé ici. Et elle est étudiée depuis les années 1950 sous le
nom de \emph{théorie de la diffusion inverse}, avec les travaux de Gelfand, Levitan et Marchenko. Ils
établissent sous quelles conditions la reconstruction est unique, et par quel algorithme la mener.

C'est ce corpus qui peut être importé : pas l'interprétation quantique, mais l'appareil mathématique de
reconstruction et de conditionnement.

\newpage
\subsection{Correspondance terme à terme avec Krapez}

\begin{center}
\begin{tabularx}{\textwidth}{@{}lXX@{}}
\toprule
\textbf{Objet} & \textbf{En mécanique quantique} & \textbf{Chez Krapez} \\
\midrule
$\psi$ & Amplitude de probabilité ; son module au carré est une densité de probabilité de présence
       & Intermédiaire de calcul construit à partir de la température ; aucune interprétation physique \\
\addlinespace
$V$ & Énergie potentielle, en joules ; le paysage dans lequel la particule évolue
    & Courbure relative de la racine de l'effusivité, en $\mathrm{s^{-1}}$ ; une propriété du matériau \\
\addlinespace
$E$ & Énergie de la particule ; un réel, positif ou négatif
    & $-p$, donc $-i\omega$ en régime modulé ; un complexe, sur l'axe imaginaire \\
\addlinespace
$x$ & Position dans l'espace, en mètres
    & $\xi$, profondeur comptée en racine de temps de diffusion \\
\addlinespace
Contrainte & Carré intégrable ; la probabilité totale vaut 1 & Aucune \\
\addlinespace
Question posée & Pour quelles valeurs de $E$ existe-t-il une solution acceptable ?
               & Quelle est la solution, pour $p$ complexe quelconque imposé ? \\
\bottomrule
\end{tabularx}
\end{center}

\medskip
Ce tableau dit que les deux équations ont la même forme et rien d'autre. C'est ce que Krapez appelle
une similarité formelle, et c'est pourquoi il écarte explicitement toute recherche de valeurs propres
ou d'états liés.

\subsection{Ce que l'analogie autorise, et ce qu'elle interdit}

\textbf{Ce qu'elle autorise.} Importer les méthodes de construction de solutions développées pour
l'équation de Schrödinger : potentiels solvables, transformations de Darboux, supersymétrie. C'est ce
que fait Krapez, et c'est parfaitement légitime, parce que ces méthodes ne dépendent que de la
\emph{forme} de l'équation.

Importer, sous condition, l'appareil de la théorie de la diffusion inverse pour parler d'unicité et de
conditionnement. Sous condition, parce qu'il faudra vérifier que les hypothèses de
Gelfand--Levitan--Marchenko s'appliquent quand $E$ n'est pas sur l'axe réel.

\textbf{Ce qu'elle interdit.} Toute affirmation portant sur la \emph{signification physique} des
objets. $\psi$ n'est pas une fonction d'onde de phonon. Il n'y a pas de niveaux d'énergie dans le
matériau. Aucune particule n'est décrite par cette équation. Le matériau ne devient pas quantique parce
que son équation ressemble à celle de Schrödinger.

\subsection{Où la vraie physique quantique intervient}

Il faut être précis, car c'est là que le sujet touche réellement au quantique, mais par un autre chemin
que l'analogie.

\textbf{Les phonons sont de vrais objets quantiques.} Ce sont les quanta de vibration du réseau
cristallin, des bosons obéissant à la statistique de Bose--Einstein, décrits par une relation de
dispersion et par l'équation de transport de Boltzmann.

\textbf{Et c'est de là que vient $\tau$.} Le temps de relaxation est un temps de collision entre
phonons ; sa valeur se calcule dans le cadre de Boltzmann, avec la séparation de Callaway entre
processus normaux et umklapp. C'est le sujet de la partie I.

\textbf{Les deux chemins ne se croisent pas.} Les phonons justifient l'existence de $\tau$ dans la loi
constitutive. L'analogie de Schrödinger fournit des outils pour analyser l'équation qui en résulte. Le
premier est de la physique, le second de la méthode. Les confondre est la faute qui coulerait
l'article.

\newpage
\section{Formulations à ne pas employer}

\begin{center}
\begin{tabularx}{\textwidth}{@{}XX@{}}
\toprule
\textbf{À éviter} & \textbf{À dire} \\
\midrule
« L'équation de la chaleur est une équation de Schrödinger »
& « Après transformation de Liouville, l'équation prend la forme d'une équation de Schrödinger
stationnaire » \\
\addlinespace
« $\psi$ décrit les phonons », « $\psi$ est la fonction d'onde du système »
& « $\psi$ est une fonction auxiliaire construite à partir de la température, sans interprétation
physique directe » \\
\addlinespace
« Les niveaux d'énergie du matériau »
& Il n'y en a pas : aucune quantification n'apparaît, faute de condition d'intégrabilité \\
\addlinespace
« Krapez n'a pas pensé à ajouter un temps de relaxation »
& « Le cadre de Krapez ne pose pas cette question ; nous en posons une autre » \\
\addlinespace
« Le temps de relaxation rend l'énergie réelle »
& « L'énergie acquiert une partie réelle et tourne vers l'axe réel ; son argument vaut $-45^{\circ}$ à
$\omega\tau = 1$ » \\
\addlinespace
« J'applique la mécanique quantique aux phonons »
& « J'importe des outils de théorie spectrale et de diffusion inverse pour analyser une équation de
même forme » \\
\addlinespace
« Les phonons et les quarks sont tous deux quantiques, donc les mêmes méthodes s'appliquent »
& Une analogie exige une identité de forme d'équation, vérifiable sur le papier ; une parenté de
vocabulaire n'en est pas une \\
\addlinespace
« La transformation de Liouville que je propose »
& « La transformation de Liouville, établie pour l'équation de la chaleur par Krapez (2018) » \\
\addlinespace
« La diffusivité disparaît »
& « La diffusivité est absorbée dans la coordonnée : elle sort de l'équation mais reste dans la
définition de l'axe » \\
\addlinespace
« Le problème inverse est mal posé »
& Préciser lequel des deux : l'unicité peut être acquise et le conditionnement mauvais ; ce sont deux
énoncés différents \\
\bottomrule
\end{tabularx}
\end{center}

\medskip
La dernière ligne mérite qu'on s'y arrête. Elle a été écrite dans le fichier 04 : un problème peut être
théoriquement injectif tout en étant pratiquement indistinguable, sous bruit et sur une bande de mesure
finie. Cette distinction doit rester nette dans tout ce qui sera rédigé --- c'est elle qui permet de
citer Krapez sans le contredire, et d'ajouter quelque chose par-dessus.

\end{document}
